% Part: first-order-logic
% Chapter: introduction
% Section: syntax

\documentclass[../../../include/open-logic-section]{subfiles}

\begin{document}

\olfileid[hi]{fol}{int}{syn}

\olsection{वाक्यविन्यास}

सबसे पहले हमें यह सटीक करना होगा कि चिह्नों की किन शृंखलाओं को प्रथम-क्रम
तर्क के !!{sentence}s की श्रेणी में रखा जाए। यह हम आगे करेंगे; अभी उदाहरणों से काम
चलाएँगे। प्रथम-क्रम तर्क के मूल निर्माण-घटक---उसकी शब्दावली---दो भागों में
बँटते हैं। पहले भाग में वे चिह्न हैं जिनसे हम विशिष्ट बातें कहते अथवा विशिष्ट
वस्तुओं को निर्दिष्ट करते हैं। वस्तुओं को निर्दिष्ट करने के लिए हम
!!{constant}s और निर्दिष्ट वस्तुओं के बारे में कुछ कहने के लिए !!{predicate}s
का उपयोग करते हैं। उदाहरणार्थ, किसी एक वस्तु को निर्दिष्ट करने के लिए हम
  $\Obj a$ को !!a{constant} मान सकते हैं और फिर !!{sentence}~$\Atom{\Obj P}{\Obj
a}$ द्वारा उसके बारे में कुछ कह सकते हैं। यदि आपके मन में ``$\Obj a$'' और
  ``$\Obj P$'' के अर्थ हैं, तो आप $\Atom{\Obj P}{\Obj a}$ को सामान्य
भाषा के किसी वाक्य की तरह पढ़ सकते हैं (और औपचारिक तर्क पहली बार सीखते समय
आपने शायद ऐसा किया भी होगा)। जब प्रथम-क्रम तर्क के !!{sentence}s के ऐसे सरल
उदाहरण मिल जाएँ, तब शब्दावली के दूसरे भाग---तार्किक चिह्नों (संयोजकों और परिमाणकों)---से
  अधिक जटिल वाक्य बनाए जा सकते हैं। उदाहरणार्थ, हम
  $(\Atom{\Obj P}{\Obj a} \land \Atom{\Obj Q}{\Obj b})$
  अथवा~$\lexists[\Obj x][\Atom{\Obj P}{\Obj x}]$ जैसे व्यंजक बना सकते हैं।

कठोर अधितार्किक अध्ययन के लिए आवश्यक अर्थविज्ञान की सटीक परिभाषाएँ और अपनी
!!{derivation} प्रणालियों के नियम देने से पहले हमें यह सटीक परिभाषित करना
होगा कि प्रथम-क्रम तर्क का !!a{sentence} किसे माना जाए। मूल विचार सरल है:
कुछ सरल !!{sentence}s के निर्माण में केवल !!{predicate}s और !!{constant}s का उपयोग होता है,
जैसे~$\Atom{\Obj P}{\Obj a}$। फिर संयोजकों और परिमाणकों की सहायता से उनसे
अधिक जटिल वाक्य बनाए जाते हैं। किंतु वे ठीक कौन-से नियम हैं जिनके अनुसार
अधिक जटिल !!{sentence}s के निर्माण की अनुमति है? इन्हें निर्दिष्ट करना आवश्यक है,
अन्यथा हमने ``प्रथम-क्रम तर्क का !!{sentence}'' पर्याप्त सटीकता से परिभाषित
नहीं किया। यहाँ कुछ समस्याएँ हैं। पहली, ठीक उन्हीं शृंखलाओं को !!{sentence}s
की श्रेणी में रखना है जिन्हें उसमें होना चाहिए। दूसरी, यह ऐसे करना है कि हम \emph{सभी}
!!{sentence}s के बारे में गणितीय प्रमाण दे सकें। अंततः !!{sentence}s पर कुछ
  आरंभिक संक्रियाओं की भी सटीक परिभाषाएँ देनी होंगी, जैसे ``प्रत्येक $\Obj x$
  को~$!A$ में~$\Obj b$ से बदलें।'' कठिनाई यह है कि प्रथम-क्रम तर्क के
परिमाणकों तथा !!{variable}s के मेल से यह काम पूरी तरह स्पष्ट नहीं रहता।
  उदाहरणार्थ, क्या $\lexists[\Obj
x][\Atom{\Obj P}{\Obj a}]$ को !!a{sentence} माना जाए?
  $\lexists[\Obj x][\lexists[\Obj x][\Atom{\Obj P}{\Obj x}]]$ का क्या? और
  ``$\Obj x$ को~$\Obj b$ से $(\Atom{\Obj
P}{\Obj x} \land \lexists[\Obj x][\Atom{\Obj P}{\Obj x}])$ में बदलें'' का
  परिणाम क्या होना चाहिए?

\end{document}
